\documentclass[11pt,a4paper]{article}
\usepackage[margin=1.8cm]{geometry}
\usepackage{tikz}
\usetikzlibrary{calc,intersections,angles,quotes,positioning}
\usepackage{multicol}
\usepackage{enumitem}
\usepackage{amsmath,amssymb}
\setlength{\parskip}{0.6em}
\setlength{\parindent}{0pt}

\definecolor{segmentblue}{RGB}{200,220,255}
\definecolor{sectororange}{RGB}{255,220,180}
\definecolor{lightgray}{RGB}{230,230,230}
\definecolor{segmentgreen}{RGB}{210,250,210}
\definecolor{segmentpurple}{RGB}{235,210,255}
\definecolor{sectoryellow}{RGB}{255,245,180}

\tikzset{
  chord/.style={thick},
  circleline/.style={thick},
  point/.style={fill=black, circle, inner sep=1.2pt},
  anglemark/.style={draw, thick},
  labelAB/.style={font=\small}
}

% Draws left pair of equal angles and, on the right, the corresponding shaded segment.
% #1 x-left, #2 y, #3 R, #4 angA, #5 angB, #6 angC, #7 angD,
% #8 x-right offset, #9 label, #10 right-start angle, #11 right-delta angle.
\newcommand{\ShadedTheoremPair}[11]{%
  \begin{scope}[shift={(#1,#2)}]
    \node[anchor=west, font=\bfseries] at (-0.2-#3, #3+0.7) {#9};
    \draw[circleline] (0,0) circle (#3);
    \coordinate (A) at ({#3*cos(#4)}, {#3*sin(#4)});
    \coordinate (B) at ({#3*cos(#5)}, {#3*sin(#5)});
    \coordinate (C) at ({#3*cos(#6)}, {#3*sin(#6)});
    \coordinate (D) at ({#3*cos(#7)}, {#3*sin(#7)});
    \draw[chord] (A)--(B);
    \draw (A)--(C)--(B);
    \draw (A)--(D)--(B);
    \pic[draw, angle eccentricity=1.2, angle radius=0.5cm] {angle=A--C--B};
    \pic[draw, angle eccentricity=1.2, angle radius=0.5cm] {angle=A--D--B};
    \fill (A) circle (1pt);
    \fill (B) circle (1pt);
    \fill (C) circle (1pt);
    \fill (D) circle (1pt);
    \node[labelAB, below right=1pt of A] {\(A\)};
    \node[labelAB, below left=1pt of B] {\(B\)};
    \node[labelAB, above right=1pt of C] {\(C\)};
    \node[labelAB, above left=1pt of D] {\(D\)};
  \end{scope}
  \begin{scope}[shift={({#1 + #8},#2)}]
    \coordinate (A2) at ({#3*cos(#4)}, {#3*sin(#4)});
    \coordinate (B2) at ({#3*cos(#5)}, {#3*sin(#5)});
    \coordinate (S) at ({#3*cos(#10)}, {#3*sin(#10)});
    \coordinate (E) at ({#3*cos(#10+#11)}, {#3*sin(#10+#11)});
    \fill[segmentblue] (S) arc[start angle=#10, delta angle=#11, radius=#3] -- (E) -- cycle;
    \draw[circleline] (0,0) circle (#3);
    \draw[very thick] (S) arc[start angle=#10, delta angle=#11, radius=#3];
    \draw[chord] (A2)--(B2);
    \fill (A2) circle (1pt);
    \fill (B2) circle (1pt);
    \node[labelAB, below right=1pt of A2] {\(A\)};
    \node[labelAB, below left=1pt of B2] {\(B\)};
  \end{scope}
}

\begin{document}

\begin{center}
  {\LARGE \bfseries Angles in the Same Segment --- Worked Solutions}\\[0.25em]
  \small Angles subtended by the same chord at the circumference, in the same segment, are equal.
\end{center}

\hrule
\vspace{1.5em}

\section*{1. Identifying the segments}

A \textbf{segment} is the region bounded by a chord and one of the arcs it cuts off.  
A \textbf{sector} is the region bounded by two radii and the arc between them.

\begin{center}
\begin{tabular}{c|c|c|c}
 & Left & Middle & Right \\ \hline
Row 1 & \(\checkmark\) segment & sector & neither \\
Row 2 & \(\checkmark\) segment & sector & \(\checkmark\) segment \\
Row 3 & \(\checkmark\) segment & sector & neither \\
Row 4 & \(\checkmark\) segment & sector & \(\checkmark\) segment
\end{tabular}
\end{center}

The shaded region in Row 2, Right is the \emph{major} segment; it is still a segment.

\section*{2. Angles in a given segment}

The vertex of each angle must lie on the circumference of the shaded segment.  
One valid set of angles is shown below. In each case the marked angles all stand on chord \(AB\).  
Any other points on the same shaded arc would also work; by the theorem all these angles are equal.

\begin{center}
\begin{tikzpicture}[scale=1]

  % (a) one angle
  \begin{scope}[shift={(-6.3,0)}]
    \coordinate (Aa) at ({2.2*cos(-50)},{2.2*sin(-50)});
    \coordinate (Ba) at ({2.2*cos(60)},{2.2*sin(60)});
    \coordinate (Ca) at ({2.2*cos(5)},{2.2*sin(5)});
    \fill[segmentblue] (Aa) arc[start angle=-50,end angle=60,radius=2.2] -- (Ba) -- cycle;
    \draw[circleline] (0,0) circle (2.2);
    \draw[chord] (Aa)--(Ba);
    \draw[thick] (Aa)--(Ca)--(Ba);
    \fill (Aa) circle (1.2pt);
    \fill (Ba) circle (1.2pt);
    \fill (Ca) circle (1.2pt);
    \node[labelAB, below right=1pt of Aa] {\(A\)};
    \node[labelAB, above right=1pt of Ba] {\(B\)};
    \node[labelAB, above right=1pt of Ca] {\(C\)};
    \pic[draw, angle radius=0.4cm, angle eccentricity=1.25] {angle=Aa--Ca--Ba};
    \node at (0,-2.9) {(a)};
  \end{scope}

  % (b) two angles
  \begin{scope}[shift={(-1.8,0)}]
    \coordinate (Ab) at ({2.2*cos(20)},{2.2*sin(20)});
    \coordinate (Bb) at ({2.2*cos(160)},{2.2*sin(160)});
    \coordinate (Cb) at ({2.2*cos(60)},{2.2*sin(60)});
    \coordinate (Db) at ({2.2*cos(120)},{2.2*sin(120)});
    \fill[segmentblue] (Ab) arc[start angle=20,end angle=160,radius=2.2] -- (Bb) -- cycle;
    \draw[circleline] (0,0) circle (2.2);
    \draw[chord] (Ab)--(Bb);
    \draw[thick] (Ab)--(Cb)--(Bb);
    \draw[thick] (Ab)--(Db)--(Bb);
    \fill (Ab) circle (1.2pt);
    \fill (Bb) circle (1.2pt);
    \fill (Cb) circle (1.2pt);
    \fill (Db) circle (1.2pt);
    \node[labelAB, below right=1pt of Ab] {\(A\)};
    \node[labelAB, below left=1pt of Bb] {\(B\)};
    \node[labelAB, above right=1pt of Cb] {\(C\)};
    \node[labelAB, above left=1pt of Db] {\(D\)};
    \pic[draw, angle radius=0.35cm, angle eccentricity=1.3] {angle=Ab--Cb--Bb};
    \pic[draw, angle radius=0.35cm, angle eccentricity=1.3] {angle=Ab--Db--Bb};
    \node at (0,-2.9) {(b)};
  \end{scope}

  % (c) three angles
  \begin{scope}[shift={(2.9,0)}]
    \coordinate (Ac) at ({2.2*cos(10)},{2.2*sin(10)});
    \coordinate (Bc) at ({2.2*cos(240)},{2.2*sin(240)});
    \coordinate (Cc) at ({2.2*cos(80)},{2.2*sin(80)});
    \coordinate (Dc) at ({2.2*cos(140)},{2.2*sin(140)});
    \coordinate (Ec) at ({2.2*cos(210)},{2.2*sin(210)});
    \fill[segmentblue] (Ac) arc[start angle=10,end angle=240,radius=2.2] -- (Bc) -- cycle;
    \draw[circleline] (0,0) circle (2.2);
    \draw[chord] (Ac)--(Bc);
    \draw[thick] (Ac)--(Cc)--(Bc);
    \draw[thick] (Ac)--(Dc)--(Bc);
    \draw[thick] (Ac)--(Ec)--(Bc);
    \fill (Ac) circle (1.2pt);
    \fill (Bc) circle (1.2pt);
    \fill (Cc) circle (1.2pt);
    \fill (Dc) circle (1.2pt);
    \fill (Ec) circle (1.2pt);
    \node[labelAB, below right=1pt of Ac] {\(A\)};
    \node[labelAB, below left=1pt of Bc] {\(B\)};
    \node[labelAB, above right=1pt of Cc] {\(C\)};
    \node[labelAB, above left=1pt of Dc] {\(D\)};
    \node[labelAB, below left=1pt of Ec] {\(E\)};
    \pic[draw, angle radius=0.35cm, angle eccentricity=1.2] {angle=Ac--Cc--Bc};
    \pic[draw, angle radius=0.35cm, angle eccentricity=1.2] {angle=Ac--Dc--Bc};
    \pic[draw, angle radius=0.35cm, angle eccentricity=1.2] {angle=Ac--Ec--Bc};
    \node at (0,-2.9) {(c)};
  \end{scope}

  % (d) ten angles
  \begin{scope}[shift={(-1.8,-6.2)}]
    \coordinate (Od) at (0,0);
    \coordinate (Ad) at ({3*cos(80)},{3*sin(80)});
    \coordinate (Bd) at ({3*cos(290)},{3*sin(290)});
    \fill[segmentblue] (Ad) arc[start angle=80,end angle=290,radius=3] -- (Bd) -- cycle;
    \draw[circleline] (0,0) circle (3);
    \draw[chord] (Ad)--(Bd);
    \fill (Ad) circle (1.5pt);
    \fill (Bd) circle (1.5pt);
    \node[labelAB, above right=1pt of Ad] {\(A\)};
    \node[labelAB, below right=1pt of Bd] {\(B\)};

    \foreach \ang/\lab in {95/$P_1$,115/$P_2$,135/$P_3$,155/$P_4$,175/$P_5$,195/$P_6$,215/$P_7$,235/$P_8$,255/$P_9$,275/$P_{10}$} {
      \coordinate (P) at ({3*cos(\ang)},{3*sin(\ang)});
      \draw[gray!60, thin] (Ad)--(P)--(Bd);
      \fill (P) circle (1.2pt);
      \node[font=\scriptsize] at ($(P)!1.18!(Od)$) {\lab};
    }

    % Thicken one representative angle.
    \coordinate (Rep) at ({3*cos(175)},{3*sin(175)});
    \draw[very thick] (Ad)--(Rep)--(Bd);
    \fill (Rep) circle (1.5pt);
    \pic[draw, angle radius=0.4cm, angle eccentricity=1.2] {angle=Ad--Rep--Bd};
    \node at (0,-3.6) {(d)};
  \end{scope}

\end{tikzpicture}
\end{center}

For (d), the grey lines show the other nine angles; the thicker angle is one of them. All ten angles are equal because they are in the same segment.

\newpage

\section*{3. Interpreting the theorem}

In each pair below, the left diagram shows the two equal angles.  
The right diagram has the correct \textbf{segment} shaded, and the heavy arc is the arc containing the vertices of the equal angles.

\begin{center}
\begin{tikzpicture}[scale=1]
  \ShadedTheoremPair{-6.8}{0}{2.2}{-25}{45}{70}{95}{6.6}{(a)}{45}{290}
  \ShadedTheoremPair{-6.8}{-5.2}{2.2}{200}{320}{145}{350}{6.6}{(b)}{320}{240}
  \ShadedTheoremPair{-6.8}{-10.4}{2.2}{150}{20}{120}{60}{6.6}{(c)}{20}{130}
  \ShadedTheoremPair{-6.8}{-15.6}{2.2}{75}{125}{170}{215}{6.6}{(d)}{125}{310}
\end{tikzpicture}
\end{center}

In (a), (b) and (d), the vertices \(C\) and \(D\) lie on the \emph{major} arc, so the major segment is shaded.  
In (c), \(C\) and \(D\) lie on the \emph{minor} arc, so the minor segment is shaded.  
In every case, the shaded region is the segment referred to by the theorem.

\end{document}